\section{Statement-by-statement comparison with the opened sources}

This section distinguishes an assertion printed in a source from a consequence
proved here using that source. The latter description is not a priority claim.
The accompanying \texttt{REVLOC\_source\_crosswalk.csv} has one row for every
numbered statement, including the unresolved statements. Bibliographical
opening notes specify the actual text consulted. In particular, opening a
publisher's two-page chapter preview is not recorded as opening a whole book.

\subsection{The manuscript passages}

We read the supplied French transcriptions at the manuscript-page markers
listed below. They are the attached, corrected transcriptions, not an
independent collation against every manuscript scan. Mathematical corrections
below therefore identify a problem in the supplied text without attributing
its origin to the author, editor, or transcriber.

\small
\begin{longtable}{@{}p{.20\textwidth}p{.72\textwidth}@{}}
\toprule
Manuscript location & Result and precise scope of the correction \\
\midrule
\endhead
XII, pp.~75--78; Proposition~6 &
The reverse implication in the sieve-cover axiom is an additional condition.
Under the modern relative-Theorem-A convention it is equivalent to the cofiber
criterion. A.3 supplies that criterion for the sphere-null examples; the
finite acyclic category in C.7 supplies a finite failure for $W_\infty$ by
adjoining two cone vertices. The implication actually printed in
Proposition~6 is W(7)$\Rightarrow$W(7 bis). \\
XII, p.~265, opening sentence; pp.~265--267 &
The reference back to Proposition~6 reverses that implication. The subsequent
cube proposition is conditional on the reverse axiom. Its forward, reverse,
and retract-dependent clauses do not establish that $W_\infty$ satisfies the
axiom. A.2--A.3 and F.2 retain the distinction. \\
XII, pp.~268--275; especially p.~270 and Cons~1--2 &
The unfamiliar reverse direction becomes exactly the assertion that an
aspherical homotopy cofiber detects its map. F.2 proves more: even the literal
absolute axioms plus the single reverse square force $W_{\mathbb Z}$.
Consequently its least solution is $W_{\mathbb Z}$, not $W_\infty$.
The supplied Cons~2 uses ``cart\'esien'' where the argument requires a
$W$-cocartesian square; the variance must be corrected to apply that argument.
No claim about the unprovided scan is made. \\
XVI, pp.~16$'$--27, \S3 &
The forward sieve operation holds for modern fundamental localizers;
the reverse operation characterizes the reversible subclass. Statements
about cubes must keep this hypothesis. R1 is checked at its used cofiber
interface in A.3; F.2 also covers the weaker historical convention. \\
XVI, p.~79, Theorem~5.4 &
Accessibility of $W_\infty$ is unaffected. The parenthetical assertion of all
Loc axioms is false for its reverse-cofibration clauses. B.1--B.4 give the
separate bounded construction needed here, including desuspension; G.7
settles accessibility of the full sphere-null example in ZFC. \\
XVI, pp.~83--86, Theorem~6.1, Lemma~6.2, Corollary~6.3 &
These fibration implications do not require reverse excision. Cisinski's
minimality theorem identifies the non-total modern possibility contained in
$W_\infty$ with $W_\infty$ itself. One cannot add reverse excision to that
non-total conclusion: reversible localizers with enough intrinsic fibrations
are total, as in the checked R3 interface. \\
XVI, pp.~87--89, \S6.4(b) &
The proposed full stability list, when it includes reverse excision and the
stated fibration condition, has only the total solution. It cannot have a
solution contained in $W_\infty$. Dropping the reverse clause changes this
conclusion. This is distinct from the legitimate problem of elementary
operations preserving simple homotopy type. \\
XVI, pp.~88--96, \S\S6.4--6.5 &
I.1--I.3 separate adjunctions, subdivision, and weak-point deletions. I.4--I.5
state the marked torsion invariant precisely. The hexagon/subdivision
observation is consistent with these results. No fundamental localizer can
refine $W_\infty$ to distinguish simple types, although torsion of a fixed
homotopy equivalence is homotopy invariant. I.6--I.8 show why the proposed
strict-globe and Thomason-cofibration substitutes do not repair the problem.
The supplied text explicitly marks manuscript p.~94 as missing. \\
XVI, p.~105, \S7.5 &
The existence of a localized derivator does not imply enough intrinsic
fibrations. The accessible example $W_{\mathbb Z}$ already disproves that
implication. E.4 determines its intrinsic fibrations exactly, and E.8 gives
the bounded unit-fiber obstruction to a factorization of any specified map. \\
XVI, p.~118, Corollary~7.7.2 and footnotes~55--56 &
F.1 refutes absolute Theorem A $\Rightarrow$ relative Theorem A, even with
Loc (8),(9). F.1.1 proves that adding Loc ($6'$ a) and (8) suffices for the
relative axiom; (9) is unnecessary. Thus the asserted corollary has a
noncircular repair, while the stronger unqualified implication is false.
The suggested base $[1]$ is terminal and cannot be a counterexample. \\
XVI, pp.~119--123, \S7.8 and footnotes~57--58 &
G.2 proves $W=W_{\mathrm{HOT}_W}$ for accessible modern fundamental
localizers. F.1.1 supplies the missing relative axiom under the listed
historical forward-gluing hypotheses. Cisinski's Theorem~6.1.18 rules out
the hoped-for smaller-than-$W_\infty$ weak localizer. \\
XVI, pp.~123--124, \S7.9 &
A left Bousfield model structure does not by itself give an intrinsic
fibration factorization: the missing issue is right properness. E.6 gives
the exact localization-unit criterion. F.3e explains what the ordinary
proper higher Thomason structures do give, separately from their arbitrary
left localizations. \\
\bottomrule
\end{longtable}
\normalsize

\subsection{The previously supplied returns}

POLES F.14 contains two separate questions. Its first, accessible reversible
$\Rightarrow RT$-fixed, is refuted by B.3--B.4, and even the full proposed
sphere-null class is accessible by G.7. Its second question, identifying
arbitrary spectral localizing subcategories with homological Bousfield
classes, is not a consequence of this answer. POLES C.10 gives bounds for
homological intrinsic fibrations. E.4 replaces those bounds by an exact
criterion for integral homology; E.1 explains why the shorter comma criterion
in the present prompt fails. The unresolved implication recorded as POLES
A.16 is settled by F.1 and F.1.1. Its description of $[1]$ as a nonterminal
test base is incorrect.

For GALOIS B.7, the properness/aspherical-generation interface is checked in
E.6 against Cisinski's Proposition~6.1.8 and Theorem~6.1.11.
E.7 and E.9 settle respectively the Morava-$K$ and $KU$ cases left among the
spectral examples of GALOIS B.14. They do not claim a classification for
every spectrum or settle every unenhanced-derivator question listed there.
H.1--H.4 answer the accessibility questions specifically requested here,
with the universes and the class form of Vop\v enka's principle made explicit.

The uses of PRETOT, INDLOCAL, 317 and 312 are likewise limited to checked
interfaces. C.7 rebuilds the acyclic binary-icosahedral witness and verifies
its finite category independently. H5.1 audits INDLOCAL B.3, B.4 and B.6;
I.9c audits D.16--D.17. F.3a reads and checks the higher Thomason transfer;
I.312 checks the elementary-expansion argument and supplies Henry's exact
comparison statements. The supplied prose claim that a Lean check exists
is not represented as an independently executed proof-assistant check.
No assertion of PATHS is needed.

\subsection{The localization and excision literature}

\paragraph{Bousfield's 1994 nullity classes.}
The opened paper \cite{JBousfield1994} constructs space nullification in
Theorem~2.2. Theorem~4.1 gives its mixed fiber sequence; Theorem~4.3 gives
a fiber-preservation conclusion under its additional base hypothesis.
Section~9.9 defines a stable comparison of \emph{space nullity classes} by
finite suspension. That terminology does not identify a localizing
subcategory of spectra. Theorems~9.12, 9.14 and 9.15 compare torsion
suspension spaces using stable nullity, connectivity and finite chromatic
type. Those hypotheses differ from the cofiber-defined reversible classes
here. They neither state B.3 nor identify the interval in D.3.2.
The explicit retractions and convergence argument in D.1.7 supply the
single-target Moore-space result used in this manuscript.

\paragraph{Bousfield's 1997 map localizations.}
In \cite{EH5Bousfield1997}, Theorem~2.2 concerns localization at a map,
Theorem~4.2 permits small intersections, and Theorem~4.4 supplies a
nullification with the same acyclic \emph{objects}. It does not identify
the two classes of inverted \emph{maps}; D.0 gives an explicit reason that
distinction matters here. Theorem~5.2 is the loop/localization comparison
used in E.7--E.9. Section~2.6's cohomological existence result imposes the
finiteness of the groups $\pi_*(E/p)$; it is not a blanket existence theorem
for arbitrary coorthogonality classes. The separate compactness proof in
G.7 is therefore necessary for the full sphere class.

\paragraph{Dror Farjoun 1996.}
The publisher previews actually opened are Chapter~1, pp.~1--2, and
Chapter~3, pp.~59--60 of \cite{JFarjoun1996}. Definition~1.A.1 gives
map-locality, and Theorem~3.A.1 compares $L_f\Omega X$ with
$\Omega L_{\Sigma f}X$ for connected $X$; 3.A.2 is its cellularization
counterpart. These statements concern specified map localization and
cellularization functors. They do not say that an arbitrary inverse
mapping-space null class is desuspension-closed. The proposed
infinite-loop sufficient condition is already refuted by D.1.1.
The rest of the book was not opened, so no negative claim about what it
might contain is made.

\paragraph{Berger--Ratkovi\'c 2019.}
Definition~2.1 of \cite{JBergerRatkovic2019} calls a left proper pointed model
category excisive when its derived suspension is conservative, and stably
excisive when suspension is fully faithful. Propositions~2.3--2.4 supply
the associated cofiber tests; Proposition~2.6 compares homotopy pushouts
and parallel cofibers. G.3, with existence supplied by G.7, gives an
excisive example whose suspension is neither faithful nor full.
Their stronger embedding discussion and the monad hypotheses of
Theorem~3.13 cannot convert conservativity into stable detection.
Thus the accessible counterexample is compatible with their results.

\paragraph{Casacuberta--Scevenels--Smith 2005.}
Proposition~6.2 of \cite{FGHCSS2005} is the reduced-product obstruction used
in H.3. Theorems~6.3--6.4 give non-smallly-generated group and space
localizations when measurable cardinals are absent. F.3c adds the model
localization verification required by the literal wording of this prompt;
H.3 constructs the derivator and checks its locally constant coefficient
class. H5.9 uses a different Heisenberg reduced-product argument for
finite-dimensional rational-unipotent tests. Those added conclusions are
proved here, not attributed verbatim to the 2005 theorems. The positive
accessibility statement in H.2 is checked against Cisinski's
Scholie~4.2.17 with its universe qualification.

\subsection{The set-size question that is not settled by these comparisons}

The attached Hovey--Palmieri text \cite{JHoveyPalmieri1999}, Section~9,
Conjecture~9.1 and Proposition~9.2, separates the assertion that every
localizing subcategory is a homological Bousfield class from its principal
case. Its proof uses the set of homological Bousfield classes only after
assuming that extra identification. This does not bound all localizing
subcategories unconditionally, and does not bound the fiber of $T$ on
reversible localizers.

The opened \cite{JCGR2014}, Theorem~3.7, establishes coreflectivity for
singly generated semilocalizing subcategories; Theorem~3.9 supplies single
generation and coreflectivity under Vop\v enka's principle. Theorem~4.2
and Corollary~4.3 concern orthogonality and nullity classes. Its discussion
on pp.~26--27 distinguishes the known proper class of one-sided nullity
classes from the set-size question with all suspension shifts. None of
those statements settles the fixed-stable-part question D.3.2. We record
what these opened theorems imply, without asserting that the article's
historical open-problem list is an exhaustive report of the literature
in November 2026.

The present return therefore does not infer a third point in the
$\loc(S/2)$ interval from any of these cardinality results. Its two
distinct endpoints are proved, and every additional strictness assertion
must supply a further unstable separating witness. The current precise
state, attempted constructions and next tests are recorded at D.3.2 and
in the research report.
